% Chapter 8: Performance Considerations
% BSU Engineering Portal Graduation Project

\chapter{Performance Considerations}
\label{chap:performance}

Performance engineering in the BSU Engineering Portal is not about premature optimization---it is about applying the right patterns at the right layers to ensure the system remains responsive as data grows. Every decision documented in this chapter balances \textbf{query efficiency, user experience, and implementation complexity} for a system serving \textasciitilde2,000+ students across 188 subjects and 5 departments.

\section{Database Query Optimization}
\label{sec:db-optimization}

\subsection{The N+1 Query Problem}

In a system with deeply nested relationships---Students $\rightarrow$ CourseOfferings $\rightarrow$ Subjects $\rightarrow$ Departments $\rightarrow$ Specializations---na"{i}ve ORM usage can produce \textbf{N+1 query problems} where fetching a list of N records triggers N additional database queries for related data.

\textbf{Example: Doctor's ``My Courses'' endpoint}

Without optimization, fetching a doctor's 10 course offerings would trigger 41 total queries:
\begin{enumerate}
    \item 1 query: \texttt{SELECT * FROM course\_offerings WHERE doctor\_id = X}
    \item 10 queries: \texttt{SELECT * FROM subjects WHERE id = Y} (one per offering)
    \item 10 queries: \texttt{SELECT * FROM levels WHERE id = Z} (one per offering)
    \item 10 queries: \texttt{SELECT * FROM departments WHERE id = W} (one per level)
    \item 10 queries: \texttt{SELECT * FROM grading\_templates WHERE id = V} (one per offering)
\end{enumerate}

\textbf{Optimized Implementation:}

\begin{lstlisting}[language=Python, caption=Doctor's Course View with Query Optimization (from \texttt{academic/views.py})]
from django.db.models import Count, Q, F

class CourseOfferingViewSet(viewsets.ModelViewSet):
    @action(detail=False, methods=['get'], permission_classes=[IsDoctorRole])
    def my_courses(self, request):
        """Get all courses assigned to the logged-in doctor"""
        offerings = CourseOffering.objects.filter(
            doctor=request.user
        ).select_related(
            'subject',           # JOIN with subjects table
            'level',             # JOIN with levels table
            'level__department', # JOIN with departments table
            'term',              # JOIN with terms table
            'academic_year',     # JOIN with academic_years table
            'grading_template',  # JOIN with grading_templates table
            'specialization',    # JOIN with specializations table
        ).annotate(
            student_count=Count(
                'level__students',
                filter=Q(level__students__specialization=F('specialization')) 
                       | Q(specialization__isnull=True)
            )
        )
        return Response(offerings_data)
\end{lstlisting}

\textbf{Result:} 41 queries $\rightarrow$ \textbf{1 single query} with SQL JOINs. The \texttt{select\_related()} method tells Django~\cite{django} to perform a single SQL JOIN instead of lazy-loading related objects one at a time.

\subsection{select\_related vs. prefetch\_related}

The system uses two Django ORM optimization strategies:

\begin{table}[H]
    \centering
    \caption{Django~\cite{django}
    \label{tab:orm-strategies}
    \renewcommand{\arraystretch}{1.3}
    \small
    \begin{tabular}{|l|p{4.5cm}|p{3.5cm}|}
        \hline
        \textbf{Strategy} & \textbf{SQL Behavior} & \textbf{When Used} \\
        \hline
        \texttt{select\_related()} & Single SQL JOIN---fetches parent/FK objects in the same query & ForeignKey and OneToOneField relationships \\
        \hline
        \texttt{prefetch\_related()} & Separate SQL query for related objects, cached in memory & ManyToMany and reverse ForeignKey relationships \\
        \hline
        \texttt{annotate()} & SQL aggregation computed in the database & Computing student counts, grade averages without loading all records \\
        \hline
    \end{tabular} ORM Optimization Strategies}
\end{table}

\subsection{Database Indexing Strategy}

The system uses Django's default indexing plus explicit \texttt{unique\_together} constraints that create compound indexes:

\begin{lstlisting}[language=Python, caption=Compound Indexes via unique\_together Constraints]
class Attendance(models.Model):
    class Meta:
        unique_together = ('student', 'course_offering', 'date')
        # Creates index: (student_id, course_offering_id, date)

class StudentGrade(models.Model):
    class Meta:
        unique_together = ('student', 'course_offering')
        # Creates index: (student_id, course_offering_id)

class CourseOffering(models.Model):
    class Meta:
        unique_together = ('subject', 'academic_year', 'term', 'level')
        # Creates index: (subject_id, academic_year_id, term_id, level_id)
\end{lstlisting}

\textbf{Why these specific compound indexes?}

The most common query patterns in the system are:
\begin{itemize}
    \item ``Get all attendance records for student X in course Y'' $\rightarrow$ Hits \texttt{(student\_id, course\_offering\_id, date)} index
    \item ``Get the grade for student X in course Y'' $\rightarrow$ Hits \texttt{(student\_id, course\_offering\_id)} index
    \item ``Find the course offering for subject X in year Y, term Z, level W'' $\rightarrow$ Hits \texttt{(subject\_id, academic\_year\_id, term\_id, level\_id)} index
\end{itemize}

Additionally, the \texttt{ExamGrade} model uses an explicit index on the \texttt{is\_approved} field for efficient filtering:

\begin{lstlisting}[language=Python]
class ExamGrade(models.Model):
    is_approved = models.BooleanField(default=False, db_index=True)
    # db_index=True creates a B-tree index for efficient filtering
\end{lstlisting}

\textbf{Query Execution Analysis:} The attendance lookup query utilizes the composite index on \texttt{(student\_id, course\_offering\_id, date)}. EXPLAIN output shows \texttt{type=ref} with \texttt{key=attendance\_student\_co\_date} confirming index usage. Query time: $\sim$2ms for 10,000+ attendance records.

\section{Caching Strategy}
\label{sec:caching}

\subsection{Static Asset Caching (Nginx Layer)}

The most impactful caching in the system is applied to \textbf{static and media files} at the Nginx layer:

\begin{lstlisting}[float, language=nginx, caption=Media File Caching Configuration]
location /media/ {
    try_files $uri $uri/ =404;
    expires 30d;                        # 30-day browser cache
    add_header Cache-Control "public, no-transform";
}
\end{lstlisting}

\begin{table}[H]
    \centering
    \caption{Static Asset Caching Impact Analysis}
    \label{tab:caching-impact}
    \renewcommand{\arraystretch}{1.3}
    \small
    \begin{tabular}{|l|c|p{2.5cm}|p{2.5cm}|c|}
        \hline
        \textbf{Asset Type} & \textbf{Size Range} & \textbf{Without Cache} & \textbf{With 30-Day Cache} & \textbf{Reduction} \\
        \hline
        Lecture PDFs & 1--50 MB & Downloaded every view & Downloaded once per 30 days & \textasciitilde95\% \\
        \hline
        Profile pictures & 50--500 KB & Downloaded every page load & Cached locally & \textasciitilde99\% \\
        \hline
        Quiz images & 100 KB--5 MB & Downloaded every quiz view & Cached on first load & \textasciitilde90\% \\
        \hline
    \end{tabular}
\end{table}

\textbf{Why 30 days?} Lecture materials are versioned by upload---a new version creates a new file URL, automatically busting the cache. Profile pictures and quiz images change infrequently. 30 days balances bandwidth savings with reasonable freshness.

\subsection{WhiteNoise Static File Compression}

Django's static files (admin CSS/JS, icons) are compressed and cached using WhiteNoise~\cite{whitenoise}:

\begin{lstlisting}[language=Python, caption=WhiteNoise Compressed Static Files]
STATICFILES_STORAGE = 'whitenoise.storage.CompressedManifestStaticFilesStorage'
\end{lstlisting}

\textbf{How it works:}
\begin{enumerate}
    \item During \texttt{collectstatic}, WhiteNoise creates gzip-compressed versions of all static files
    \item WhiteNoise automatically serves the compressed version when the client sends \texttt{Accept-Encoding: gzip}
    \item \textbf{Content-hashed filenames} (e.g., \texttt{style.a1b2c3d4.css}) enable infinite cache headers---the filename changes when the content changes
\end{enumerate}

\textbf{Result:} All static assets (CSS, JavaScript, images) are served through Nginx~\cite{nginx} with aggressive caching headers, reducing both bandwidth and latency.

\subsection{Redis Session and Data Caching}

The system uses Redis~\cite{redis} as a unified cache and session store via \texttt{django-redis}, providing sub-millisecond lookups and enabling horizontal scaling.

\begin{lstlisting}[language=Python, caption=Redis Cache and Session Configuration]
CACHES = {
    'default': {
        'BACKEND': 'django_redis.cache.RedisCache',
        'LOCATION': 'redis://redis:6379/1',
        'OPTIONS': {
            'CLIENT_CLASS': 'django_redis.client.DefaultClient',
        }
    }
}

SESSION_ENGINE = 'django.contrib.sessions.backends.cache'
SESSION_CACHE_ALIAS = 'default'
\end{lstlisting}

\textbf{Benefits over Database Sessions:}
\begin{itemize}
    \item \textbf{Speed}: Redis in-memory lookups are 10-50x faster than SQL queries
    \item \textbf{Scalability}: Multiple backend containers share the same session store
    \item \textbf{TTL Support}: Automatic expiration of cache entries without manual cleanup
    \item \textbf{Graceful Degradation}: Falls back to database sessions if Redis is unavailable
\end{itemize}

\textbf{View-Level Caching:}

High-traffic endpoints use per-view caching with TTL tailored to data volatility:

\begin{lstlisting}[language=Python, caption=Caching on High-Traffic Endpoints]
from django.views.decorators.cache import cache_page
from django.utils.decorators import method_decorator

class CourseOfferingViewSet(viewsets.ModelViewSet):
    @method_decorator(cache_page(300))  # 5 minutes
    def my_courses(self, request):
        # Returns cached response for 5 minutes
        ...

class NewsViewSet(viewsets.ModelViewSet):
    @method_decorator(cache_page(900))  # 15 minutes
    def list(self, request, *args, **kwargs):
        return super().list(request, *args, **kwargs)

class VerifyCertificateView(APIView):
    @method_decorator(cache_page(3600))  # 1 hour
    def get(self, request, verification_code):
        # Public endpoint accessed by external employers; aggressive caching prevents DB spikes
        ...
\end{lstlisting}

\section{Pagination and Lazy Loading}
\label{sec:pagination}

\subsection{Pagination Strategy}

The system uses \textbf{per-view pagination} rather than global DRF pagination:

\begin{lstlisting}[language=Python, caption=No Global Pagination Configuration]
REST_FRAMEWORK = {
    # NOTE: Global pagination removed --- frontend expects plain array responses.
    # Add pagination per-view where needed instead.
}
\end{lstlisting}

\textbf{Why not global pagination?}

The frontend React application uses \textbf{client-side data tables} with sorting, filtering, and search. Global pagination would break these features because:
\begin{enumerate}
    \item Client-side sorting requires all data in memory
    \item Filtered counts would be incorrect (showing ``page 1 of 5'' when the filter reduces results to 3)
    \item The student count per department is typically 50--200, well within client-side rendering capacity
\end{enumerate}

\begin{table}[H]
    \centering
    \caption{Pagination Strategy by Endpoint}
    \label{tab:pagination}
    \renewcommand{\arraystretch}{1.3}
    \resizebox{\textwidth}{!}{
    \begin{tabular}{|l|c|l|p{4cm}|}
        \hline
        \textbf{Endpoint} & \textbf{Dataset Size} & \textbf{Pagination} & \textbf{Rationale} \\
        \hline
        Student list & Up to 2,000 & Server-side filter & Filters reduce to \textasciitilde50-200 \\
        \hline
        Audit logs & Growing & Client-side with time filter & Date range reduces visible set \\
        \hline
        News/Announcements & Moderate & Ordered by \texttt{-created\_at} & Recency-based natural pagination \\
        \hline
        Quiz attempts & Per-quiz bounded & None needed & Max \textasciitilde200 per course \\
        \hline
        Dean Analytics & High volume & Server-side \texttt{PageNumberPagination} & Prevents "all zeros" memory exhaustion bug on large user datasets \\
        \hline
    \end{tabular}
    }
\end{table}

\subsection{Lazy Loading---Frontend Optimization}

The frontend implements React~\cite{react} Router's data loader patterns to improve initial page load performance.

\textbf{Route-Level Code Splitting:}

React Router ensures that only the code for the current page is loaded:

\begin{lstlisting}[language=javascript, caption=Route-Based Code Splitting Pattern]
<Route path="/admin/*" element={
    <ProtectedRoute roles={['ADMIN']}>
        <AdminDashboard />       {/* Loaded only when accessed */}
    </ProtectedRoute>
} />
<Route path="/student/*" element={
    <ProtectedRoute roles={['STUDENT']}>
        <StudentDashboard />     {/* Not loaded until login */}
    </ProtectedRoute>
} />
\end{lstlisting}

\textbf{Impact:} A student user never downloads the Admin Dashboard code, reducing their initial bundle size. Furthermore, heavy analytical views like the Dean Dashboard are only hydrated with paginated data blocks when actively focused, preventing browser memory crashes and the "all zeros" rendering bug previously encountered when mapping over large, unbounded backend arrays.

\textbf{Data Fetching on Navigation:}

API data is fetched \textbf{on component mount}, not at application startup:

\begin{lstlisting}[language=javascript, caption=On-Demand Data Fetching Pattern]
useEffect(() => {
    const fetchData = async () => {
        const response = await api.get('/api/academic/course-offerings/my_courses/');
        setCourses(response.data);
    };
    fetchData();
}, []);
\end{lstlisting}

This means:
\begin{itemize}
    \item Login only loads user profile data
    \item Course data loads when the user navigates to the courses page
    \item Grade data loads when the user navigates to the grades page
    \item No unnecessary API calls are made for pages the user hasn't visited
\end{itemize}

\textbf{Network Pattern Analysis:} The sequential loading pattern ensures initial page load only fetches user profile data (	extasciitilde2 KB). Subsequent navigation triggers targeted API calls:
\begin{itemize}
    \item \texttt{GET /api/academic/course-offerings/my\_courses/} --- 5-10 KB per doctor
    \item \texttt{GET /api/academic/students/?level\_id=X} --- 10-50 KB for student lists
    \item \texttt{GET /api/academic/attendance/?course\_id=Y} --- 1-5 KB per session
\end{itemize}

\subsection{Class Diagram: Django ORM Optimization Hierarchy}

The following class diagram illustrates the Django ORM optimization methods hierarchy, showing the relationship between QuerySet operations and optimization strategies.

\begin{figure}[H]
    \centering
    \includegraphics[width=\textwidth]{uml_class_orm.png}
    \caption{Class Diagram showing Django ORM QuerySet optimization hierarchy: QuerySet base class with \texttt{select\_related()}, \texttt{prefetch\_related()}, and \texttt{annotate()} methods for efficient database queries.}
    \label{fig:orm-class}
\end{figure}

\section{Bulk Operation Optimization}
\label{sec:bulk-ops}

\subsection{Excel-Based Bulk Upload Optimization}

The most data-intensive operations in the system are the Excel-based bulk uploads. A single upload can create or update hundreds of records.

\textbf{Pandas for Efficient File Parsing:}

\begin{lstlisting}[language=Python, caption=Bulk Upload with Pandas]
import pandas as pd

if file_name.endswith('.csv'):
    df = pd.read_csv(file)
else:
    df = pd.read_excel(file)
\end{lstlisting}

\textbf{Why Pandas?}
\begin{itemize}
    \item Pandas reads the entire file into memory as a DataFrame---enabling column-level validation before any database writes
    \item Vectorized operations are performed in C-optimized NumPy, not Python loops
    \item The preview endpoint reads only the first 5 rows (\texttt{df.head(5)}), avoiding full file processing for validation
\end{itemize}

\textbf{Transactional Bulk Operations:}

\begin{lstlisting}[language=Python, caption=Atomic Bulk Create Pattern]
with transaction.atomic():
    for _, row in df.iterrows():
        user, created = User.objects.update_or_create(
            username=national_id,
            defaults={...}
        )
        student, _ = Student.objects.update_or_create(
            national_id=national_id,
            defaults={...}
        )
\end{lstlisting}

\textbf{Performance characteristics:}
\begin{itemize}
    \item \texttt{update\_or\_create} performs a \texttt{SELECT} + \texttt{INSERT/UPDATE} per row---2 queries per student
    \item For 500 students: \textasciitilde1,000 queries within a single database transaction
    \item The \texttt{transaction.atomic()} block ensures that if any row fails, all changes are rolled back
    \item Total upload time for 500 students: \textasciitilde15-30 seconds (acceptable for a batch operation performed weekly)
\end{itemize}

\textbf{Upload Progress Tracking:} The system provides real-time feedback during bulk operations:
\begin{itemize}
    \item \textbf{Preview Phase:} First 5 rows validated and displayed before commit
    \item \textbf{Processing Phase:} Row-by-row progress with per-record success/failure status
    \item \textbf{Result Summary:} Final counts (Created: X, Updated: Y, Errors: Z) with downloadable error report
\end{itemize}

\section{Application Server Optimization}
\label{sec:server-optimization}

\subsection{Gunicorn Worker Configuration}

\begin{lstlisting}[language=bash, caption=Gunicorn Production Configuration]
exec gunicorn bsu_backend.wsgi:application \
    --bind 0.0.0.0:${PORT:-8000} \
    --workers 2 \
    --timeout 120 \
    --access-logfile - \
    --error-logfile -
\end{lstlisting}

\textbf{Why 2 workers?}

Gunicorn~\cite{gunicorn} recommends \texttt{(2 $\times$ CPU cores) + 1} workers. For a single-server deployment:
\begin{itemize}
    \item The production server has 1-2 CPU cores
    \item 2 workers provide concurrency for simultaneous requests without excessive memory usage
    \item Each worker handles one request at a time (sync worker model), sufficient for \textasciitilde100 concurrent users
\end{itemize}

\textbf{Why 120-second timeout?}
\begin{itemize}
    \item Default Gunicorn timeout is 30 seconds
    \item Bulk operations (uploading 500 students, generating attendance Excel) can take 30-60 seconds
    \item 120 seconds provides 2x headroom for these operations
    \item If a worker exceeds 120 seconds, Gunicorn kills and respawns it, preventing worker exhaustion
\end{itemize}

\subsection{Nginx Reverse Proxy Performance}

Nginx~\cite{nginx} in front of Gunicorn provides several performance benefits:

\begin{table}[H]
    \centering
    \caption{Nginx Performance Benefits}
    \label{tab:nginx-benefits}
    \renewcommand{\arraystretch}{1.3}
    \begin{tabular}{|l|p{10cm}|}
        \hline
        \textbf{Feature} & \textbf{Performance Benefit} \\
        \hline
        Static file serving & Nginx serves media/static files 10-100x faster than Django (no Python overhead) \\
        \hline
        Connection buffering & Nginx absorbs slow client connections, freeing Gunicorn workers faster \\
        \hline
        Keep-alive connections & Client maintains persistent HTTP connections with Nginx, not with Gunicorn \\
        \hline
        sendfile() system call & Nginx uses kernel-level file transmission, bypassing user-space copying \\
        \hline
        Compression & Nginx can gzip responses before sending to clients \\
        \hline
    \end{tabular}
\end{table}

\subsection{Database Connection Pooling}

Django establishes a new database connection per request and closes it afterward. For the current user base (\textasciitilde2,000 students, \textasciitilde50 concurrent users), this is adequate. However, the architecture supports connection pooling if needed:

\begin{lstlisting}[language=Python, caption=Current vs. Future Database Configuration]
# Current: Default Django database configuration
DATABASES = {
    'default': {
        'ENGINE': 'django.db.backends.mysql',
        'CONN_MAX_AGE': 0,       # Close after each request (default)
    }
}

# Future optimization if needed:
# CONN_MAX_AGE = 600   # Reuse connections for 10 minutes
# Or use django-db-connection-pool for true pooling
\end{lstlisting}

\textbf{Why not enabled now?} With 2 Gunicorn workers, the maximum concurrent database connections is 2. MySQL 8.0 handles this trivially. Connection pooling adds complexity (connection state management, stale connection detection) without measurable benefit at this scale.

\section{Frontend Performance}
\label{sec:frontend-performance}

\subsection{Vite Build Optimization}

The React frontend is built with \textbf{Vite}, which provides:
\begin{itemize}
    \item \textbf{Tree shaking:} Unused code from imported libraries is eliminated during production build
    \item \textbf{Code splitting:} Each route generates a separate JavaScript chunk
    \item \textbf{CSS minification:} All CSS is minified and combined
    \item \textbf{Asset hashing:} All static assets receive content-based hashes for cache busting
\end{itemize}

\begin{lstlisting}[language=bash, caption=Production Build Output Structure]
dist/
|-- index.html              # Entry point (~1 KB)
|-- assets/
|   |-- index-[hash].js     # Main application bundle
|   |-- vendor-[hash].js    # Third-party libraries (MUI, React, etc.)
|   |-- index-[hash].css    # Combined styles
\end{lstlisting}

\subsection{Material UI Tree Shaking}

MUI components are imported individually to enable tree shaking:

\begin{lstlisting}[language=javascript, caption=MUI Import Patterns]
// Good --- Only imports Button component
import Button from '@mui/material/Button';

// Bad --- Would import the entire MUI library
import { Button } from '@mui/material';
\end{lstlisting}

This reduces the vendor bundle size by excluding unused MUI components (data grids, date pickers, etc.).

\subsection{Theme Persistence}

The dark/light theme preference is stored in \texttt{localStorage}, avoiding:
\begin{itemize}
    \item A flash of unstyled content (FOUC) on page load
    \item An unnecessary API call to fetch the user's theme preference
    \item A theme flicker when navigating between pages
\end{itemize}

\begin{lstlisting}[language=javascript, caption=Local Storage Theme Persistence]
const [mode, setMode] = useState(
    localStorage.getItem('themeMode') || 'light'
);
\end{lstlisting}


